\section{Data Engineering and Knowledge-Graph Construction}\label{sec:chapter-3-data-engineering-and-knowledge}

Hard verification is the organising principle of this dissertation, and it is only possible over data whose semantics are explicit. A verifier cannot check that a sum is legitimate unless the data records which values may be summed. An oracle audit cannot re-derive an answer unless every convention behind the original computation is written down as data rather than remembered informally. This chapter describes how we turned five years of UK Open Contracting Data Standard (OCDS) releases into such a substrate: a typed, convention-explicit, audit-first knowledge graph of 215,221 contract-award nodes and 131,502 canonical organisations [ARTIFACT-VERIFIED: data/kg/nodes/contract\_nodes.parquet, data/kg/nodes/org\_nodes.parquet]. Two design stances run through the chapter. The first is \emph{precision-first entity resolution}. A false merge corrupts every downstream count and sum, so merging is tiered, conservative, and audited. The documented cost is under-merging, not silent over-merging. The second is \emph{conventions as data}. The additivity of money values, the flat first-party record universe, and the latest-release snapshot are checkable flags and stated choices, not implicit assumptions. Chapter 4's independent oracle audit and Chapter 5's verifier both depend on exactly these properties.

\section{Source data and ingestion}\label{sec:source-data-and-ingestion}

\subsection{The OCDS release corpus}\label{sec:the-ocds-release-corpus}

The source data is the UK's public OCDS publication covering Contracts Finder and Find a Tender notices. It is distributed as one compressed JSON-Lines file per year for 2022--2026. Each line is a \emph{compiled release}: a self-contained JSON snapshot of one contracting process, identified by a stable \texttt{ocid}. A release carries the buyer reference and a \texttt{parties} array with the full organisation records. It also carries a \texttt{tender} block (title, CPV classification, estimated value, procurement method, lots, periods), an \texttt{awards} array (suppliers, award values, related lots), and a \texttt{contracts} array (signed values and dates). The five year-files contain 166,277 release records in total (source analysis: ocds\_data\_analysis.md). Volume jumps sharply in 2025, when the Procurement Act 2023 introduced new mandatory notice types. A release is a snapshot of a \emph{process}, not a contract: several releases of the same OCID can appear within and across year files as the process moves from notice to award to signature.

Ingestion (\texttt{src/procurement\_graph/ingest/loader.py}) streams each gzipped file line by line, so no file is ever loaded whole. Each release is parsed defensively: malformed lines are logged and skipped, and rows missing \texttt{ocid} or \texttt{date} are dropped. Each release is then flattened into one row, with the nested \texttt{parties} and \texttt{contracts} arrays serialised as JSON strings, so the interim Parquet file has a simple and universally readable schema. Deduplication keeps, for every OCID, the release with the latest ISO-8601 \texttt{date}; string comparison is safe because all dates are ISO-formatted. A \texttt{source\_years} column records every year-file in which the OCID was observed.

The \emph{latest-release snapshot} is the first explicit convention of this chapter, and it is a choice with alternatives. Modelling full amendment chains would preserve every correction, but it multiplies schema complexity for question families the benchmark never asks. Keeping the first release would discard corrections. Taking the latest release matches the publisher's own compiled presentation and gives exactly one authoritative snapshot per contracting process. The convention is stated here and stress-tested by the oracle audit in Chapter 4. The knowledge graph therefore answers questions about the latest published state of each process within the ingested corpus. It does not answer questions about publication history, and it does not cover UK procurement in its entirety.

A second, wider extraction pass (\texttt{src/procurement\_graph/extract/tables.py}) re-reads the raw files to capture fields the first pass does not. It produces seven tables (tender core, lots, award criteria, awards, bid statistics, text evidence, documents). Alignment with the deduplicated snapshot is enforced structurally, not by convention: a raw release is accepted only when both its \texttt{ocid} \emph{and} its release \texttt{id} match the winning release chosen at ingestion. This guarantees that the extracted tables are row-for-row consistent with the interim snapshot.

\subsection{Quality problems the pipeline must absorb}\label{sec:quality-problems-the-pipeline-must-absor}

Three families of data-quality problems shaped the design. All were documented at analysis time, and all are visible in the code paths that handle them.

\textbf{Money semantics.} Where a contract's value lives is inconsistent across the corpus. In the 2024 population, \texttt{awards[].value} was null in every inspected record. The reliable signed value sits in \texttt{contracts[].value.amount}, joined to its award through \texttt{awardID}. \texttt{tender.value.amount} is a pre-award \emph{estimate}, present in roughly a third of records (source analysis: ocds\_data\_analysis.md). Worse, the corpus mixes genuinely additive spend with framework-agreement ceilings --- multi-year umbrella values that reach tens of billions of pounds --- so a naive SUM over "contract values" is not total spend. The extraction layer therefore records a \texttt{value\_source} per award with a fixed precedence (\texttt{award}, then \texttt{contract}, then \texttt{tender}) and never overwrites provenance. The graph derives a boolean \texttt{value\_is\_additive} that is true exactly when the value came from an award or a signed contract, and false when it is a tender estimate. In the final graph, 160,362 award rows carry contract-sourced values, 15,640 award-sourced, 19,623 tender-sourced, and 19,596 no value at all; 176,002 of 215,221 rows are flagged additive [ARTIFACT-VERIFIED: data/kg/nodes/contract\_nodes.parquet]. This is the "conventions as data" stance in miniature. Instead of the pipeline privately knowing that tender estimates must not be summed, the flag turns \emph{additive-only money aggregation} into a guard that the downstream executor and verifier enforce mechanically (Chapter 5). The build itself also validates the flag: the KG validation suite fails the build if any award- or contract-sourced value is marked non-additive, or any tender fallback is marked additive (\texttt{src/procurement\_graph/kg/validate.py}).

\textbf{Organisation identity noise.} The record-level \texttt{buyer} field is a bare \texttt{{id, name}} reference. The authoritative organisation records live in \texttt{parties}, and the identifiers they carry are of very unequal quality. Around nine in ten buyer references and 86\% of supplier references use \texttt{GB-FTS} platform identifiers. These are per-profile account IDs on the publishing platform, not canonical entity identifiers: the Ministry of Defence appeared under 77 distinct GB-FTS IDs in a 2024 sample (source analysis: ocds\_data\_analysis.md). Name strings are noisy even when official identifiers exist (casing variants, ampersand versus "and", trailing whitespace, legal-name drift after mergers). Some source strings are simply malformed. Organisation names with duplicated suffixes, such as "\ldots{} NHS Trust Trust Headquarters", appear verbatim in the data and later surface in Chapter 7's error analysis. On such an unresolvable mention, the correct system behaviour is a provenance-gated abstention rather than a guessed match [WORKLOG-SOURCED: cicada\_worklog.md, 2026-07-05]. Section 3.3 describes the entity-resolution design these problems dictate.

\textbf{Structural multiplicity.} One release may hold several awards. One award may name several suppliers (\texttt{awards[].suppliers} is a list). Lots relate awards to descriptions many-to-many. The pipeline never collapses these silently: awards are extracted one row per \texttt{(ocid, award\_id)} with the supplier list serialised intact, and the graph gives multi-supplier awards one edge per supplier organisation (Section 3.2).

\section{Graph schema and coverage}\label{sec:graph-schema-and-coverage}

The knowledge graph (built by \texttt{src/procurement\_graph/kg/}, orchestrated from \texttt{pipelines/40\_build\_kg.py}) is a set of typed Parquet tables --- four node tables and four edge tables --- whose column contracts are frozen in code (\texttt{kg/schema.py}).

\begin{table}[t]
\centering
\small
\caption{Entity-resolution tiers: alias counts merged at each confidence tier, with the residual singletons.}
\label{tab:03_data_kg_1}
\begin{tabular}{lll}
\toprule
Table & Rows & Key \\
\midrule
\texttt{contract\_nodes} & 215,221 & \texttt{contract\_node\_id} = \texttt{contract:{ocid}:{award\_id}} \\
\texttt{org\_nodes} & 131,502 & \texttt{canonical\_id} \\
\texttt{cpv\_nodes} & 3,870 & \texttt{cpv\_code} \\
\texttt{evidence\_nodes} & 535,731 & \texttt{evidence\_id} \\
\texttt{buyer\_of} & 215,218 & one edge per contract node with a resolvable buyer \\
\texttt{supplier\_of} & 334,063 & one edge per (supplier organisation, contract node) \\
\texttt{categorized\_by} & 164,691 & contract $\rightarrow$ CPV code \\
\texttt{evidence\_for} & 1,326,240 & evidence $\rightarrow$ contract, with match scope \\
\bottomrule
\end{tabular}
\end{table}


All row counts verified directly [ARTIFACT-VERIFIED: data/kg/nodes/\textbackslash{}\emph{.parquet, data/kg/edges/\textbackslash{}}.parquet].

The unit of analysis is the \emph{contract-award}: one node per \texttt{(ocid, award\_id)} pair, spanning 118,088 distinct OCIDs [ARTIFACT-VERIFIED: data/kg/nodes/contract\_nodes.parquet]. Each contract node carries three groups of fields. The award's own fields cover status, title, dates, and value with \texttt{value\_source} and \texttt{value\_is\_additive}. The release-level tender context, joined in by OCID, covers title, CPV code and description, procurement method and category, and periods. Derived conveniences are computed once at build time rather than at query time: \texttt{has\_award\_signed\_date}, \texttt{has\_contract\_period}, day-count intervals, and contract duration. Organisation nodes carry the canonical name and identifiers, the entity-resolution status and alias counts, region evidence aggregated across all observations of the entity, and precomputed role summaries. The role summaries include buyer/supplier contract counts, first and last activity years, and additive-only value sums, in which non-additive values contribute zero by construction (\texttt{kg/enrichment.py}). CPV nodes carry the code, its description, its division/group/class prefixes (which make the CPV hierarchy navigable by prefix), and the same additive-only aggregates.

Coverage is reported as measured fact, not assumption. 215,218 of 215,221 contract nodes have a buyer edge. 204,186 have at least one supplier edge (334,063 supplier edges in total; 22,075 awards name more than one supplier organisation, and one names 249). 164,691 carry a CPV classification, and 215,202 have at least one text-evidence pointer [ARTIFACT-VERIFIED: data/kg/edges/\textbackslash{}*.parquet]. Among organisations, 7,125 appear as buyers, 114,281 as suppliers, and 1,269 in both roles [ARTIFACT-VERIFIED: data/kg/nodes/org\_nodes.parquet]. These numbers matter downstream because they fix abstention semantics. A question about the supplier of an award with no supplier edge must produce an empty result and an abstention, not an invented supplier. The benchmark's empty-result traps (Chapter 4) are sampled with these bounds in mind.

Why this shape? The workload is aggregate and multi-hop question answering: counts, sums, rankings, extremes, and bridge joins such as "contracts whose supplier is also a buyer" or "top suppliers within a CPV category for a given buyer". Two properties of the schema serve exactly this workload. First, every queryable property is a \emph{typed slot with a closed value space}. Procurement category is one of goods, services, or works; CPV codes come from the 3,870-row node table; organisation references resolve against the canonical table. This gives the planner in Chapter 5 concrete grounding targets, and it gives the verifier closed enums to check against. Second, the graph flattens deterministically into a single record universe. The query backend (\texttt{src/procurement\_graph/qa/benchmark/kg\_interface.py}) joins the edges into one row per contract node. Each row carries the full tuples of buyer and supplier canonical IDs, names, and regions, together with scalar \texttt{buyer\_name}/\texttt{supplier\_name} columns defined as the first element of the sorted name tuple. This \emph{flat first-party record universe} is the second stated convention. Filters and aggregations are defined over one row per contract-award with a deterministic first-party scalar. Multi-party membership remains queryable through the tuple columns and \texttt{supplier\_count}, and the independent oracle of Chapter 4 recomputes answers under exactly the same stated flattening.

A deliberate non-decision completes the schema story: the graph is columnar Parquet, not a graph database. At this scale --- a few hundred thousand nodes and about two million edges --- vectorised dataframe joins comfortably outperform query-engine round-trips. The executor stays deterministic, dependency-free, and unit-testable, and the independent oracle audit can re-derive every answer with plain dataframes and no shared query engine. This is an argument about operational surface, not expressiveness: nothing in the workload requires traversals that columnar joins cannot express.

\section{Entity resolution: conservative, tiered, audited}\label{sec:entity-resolution-conservative-tiered-au}

The alias universe contains 204,711 raw organisation identifiers, 76.9\% of them notice-scoped GB-FTS platform IDs, resolving to 131,502 canonical organisations [ARTIFACT-VERIFIED: data/entities/alias\_map.parquet, data/entities/canonical\_orgs.parquet]. A single false merge silently corrupts every count and sum that touches the merged entity. Resolution is therefore organised as ordered tiers of strictly decreasing evidence strength, each recorded per entity in an \texttt{er\_status} field, and everything below the evidence bar is left unmerged. The realised tier sizes are: registry-identifier deterministic merges 26,704; safe deterministic variants 6,332; normalised name + region 14,875; name-only 1,286; LLM-adjudicated borderline cases 255; curated government lookup 37; and 82,013 singletons deliberately left unmerged [ARTIFACT-VERIFIED: data/entities/canonical\_orgs.parquet, \texttt{er\_status} value counts].

\textbf{Phase 1 --- deterministic, registry-anchored} (\texttt{src/procurement\_graph/er/phase1.py}). Any party observed with an official registry scheme takes that registry identifier as its canonical ID. The schemes are Companies House (GB-COH), NHS organisation codes (GB-NHS), UK Provider Reference Numbers (GB-UKPRN), the charity registers (GB-CHC, GB-SC, GB-NIC), and the Mutuals register (GB-MPR). A fixed priority order applies when several schemes are observed for one raw ID. A single, tightly scoped cross-referencing rule then attaches platform IDs to these anchors: when, \emph{within one OCID}, a GB-FTS party and an officially identified party share the same normalised name, the GB-FTS ID becomes an alias of the official entity. The OCID scoping is the safety property. Co-occurrence of both identifiers on the same notice is direct evidence that they denote one organisation; the same name match across unrelated notices would not be. GB-FTS identifiers are never themselves canonical. They are notice-scoped platform artefacts, and treating them as entity keys causes exactly the fragmentation failure (one ministry, 77 nodes) the design exists to avoid.

\textbf{Phase 2 --- heuristic but bounded} (\texttt{src/procurement\_graph/er/phase2.py}). Entities left unresolved by Phase 1 pass through three steps. First comes a small curated lookup of well-known government bodies (37 canonical entities, e.g. \texttt{GOV-MOD} for the Ministry of Defence, which ultimately absorbs 1,222 aliases, the largest alias group in the table [ARTIFACT-VERIFIED: data/entities/alias\_map.parquet]). Second comes exact normalised-name + region grouping. The normaliser (\texttt{common/normalise.py}) uppercases, strips accents and punctuation, rewrites ampersands, and removes legal-form suffixes (LIMITED, LTD, PLC, LLP, CIC, \ldots{}); groups of more than one raw ID sharing a normalised name and region merge under a stable content-hashed \texttt{MERGED-*} identifier. Third, name-only merging is the fallback for entities with no region evidence. Everything still unmatched becomes a \emph{singleton} and keeps its own raw ID as a provisional canonical ID. Fuzzy string matching was deliberately demoted: the Jaro-Winkler candidate generator (\texttt{er/candidates.py}) is a read-only analysis tool with blocking and hard output caps. It produces human-review files and \emph{never} writes to the entity tables.

\textbf{Safe-merge extension --- audited, cluster-checked.} A later ablation framework (\texttt{src/procurement\_graph/experiments/er\_ablation.py}, applied by \texttt{scripts/apply\_safe\_er\_merges.py}) generated candidate pairs with risk flags (scheme conflicts, parent/child and subsidiary patterns, acronym ambiguity, procurement-agent patterns). The deterministic-safe subset and an LLM-adjudicated borderline queue then passed through cluster-level checks, and the surviving merge edges were applied by union-find. Canonical-ID election within a merged component prefers official registry IDs (in scheme-priority order), then curated \texttt{GOV-*} IDs, then existing \texttt{MERGED-*} IDs. The LLM tier contributes only 255 of 131,502 canonical entities, and every one carries \texttt{er\_status='llm\_safe\_merge'} so it can be inspected, discounted, or reverted wholesale. The two extension tiers account for the \texttt{deterministic\_safe\_merge} (6,332) and \texttt{llm\_safe\_merge} (255) statuses in the audit trail.

\textbf{The audit trail.} Every canonical entity --- merged or singleton --- has one row in \texttt{er\_audit.csv}. The row records its canonical ID, display name, resolution status, official scheme, alias count, and the full pipe-delimited list of raw IDs it absorbs (131,502 rows [ARTIFACT-VERIFIED: data/entities/er\_audit.csv]). Together with the \texttt{alias\_source} column of the alias map (phase1 / phase2 / safe-merge / unresolved), every merge decision in the graph can be reconstructed after the fact. The build also verifies referential closure: every alias target must exist as an organisation node, or the build fails.

\textbf{Residual ambiguity is handled at query time, conservatively.} The design makes no recall claim. Its documented cost is that decoration variants of one organisation --- "The X" versus "X", "X (BEIS)" versus "X", casing differences --- can survive as distinct canonical rows, splitting one organisation's records across variants. The runtime policy (\texttt{src/procurement\_graph/reasoning/entity\_resolution.py}) is the single shared gate through which every planner and repair path resolves an organisation mention, so no caller silently takes the top hit of a weak match. Its logic, directly from the code, is as follows. An exact record match or a $\geq$0.999 score resolves immediately. Otherwise the top candidate must score at least \texttt{min\_score = 0.85}, and if a second candidate also clears that bar, the top must lead by a margin of at least \texttt{min\_margin = 0.10}. There is one exception. When all contenders collapse to one \emph{decoration-variant equivalence key} (casefold; strip a leading "the" and a trailing parenthetical alias; unify "\&"/"and"; drop punctuation), the mention is resolved, but \emph{every} variant is surfaced. The executor then filters with an \texttt{IN} over all of them rather than undercounting on one. Genuine ambiguity, meaning distinct equivalence keys within the margin, refuses to resolve and feeds the abstention path. A minimal, evidence-driven mention-expansion table is applied only when the literal lookup finds nothing. It currently holds the single rule ICB $\rightarrow$ "Integrated Care Board", added when the audit showed the graph stores full names while questions use the abbreviation. Both mechanisms were introduced \emph{after} the Chapter 4 oracle audit identified variant splits and abbreviation mismatch as the ambiguity classes with practical impact. This follows the thesis-wide discipline: conventions and repairs are adopted in response to measured failures, then written down.

\section{Provenance: evidence nodes and why they exist}\label{sec:provenance-evidence-nodes-and-why-they-e}

The fourth node type stores no structured facts at all. The 535,731 evidence nodes hold the corpus's long free-text fields --- lot descriptions (332,802), tender descriptions (164,373), selection-criteria descriptions (21,429), review details, and performance terms [ARTIFACT-VERIFIED: data/kg/nodes/evidence\_nodes.parquet]. They are extracted verbatim with their OCDS field paths and cover 166,257 distinct OCIDs, essentially the whole release corpus. They are deliberately kept out of the structured record universe. Paragraphs are not queryable slots, and storing them as contract-node attributes would blur the line between what the executor can filter on and what it can only cite.

Their purpose is provenance. The \texttt{evidence\_for} edge table (1,326,240 edges reaching 215,202 of 215,221 contract nodes) links each text fragment to the contract-award nodes it describes. The link is at lot scope when the fragment's lot ID matches a specific award's related lots, and at OCID scope otherwise. The resolution is recorded in a \texttt{match\_scope} column (\texttt{lot}, \texttt{ocid}, or \texttt{ocid\_fallback}), so the precision of every provenance link is itself queryable. Downstream, this is what makes \emph{provenance-gated abstention} concrete rather than rhetorical. Every answer must trace to specific records, and answers are returned with evidence samples attached. When a question's constraints match no records, such as an unresolvable "NHS Trust Trust" mention or a year with no awards for the named buyer, the system returns an explicit empty-evidence abstention instead of a fluent guess. The benchmark's empty-result abstention family (Chapter 4) tests precisely this contract.

\section{Engineering properties that matter scientifically}\label{sec:engineering-properties-that-matter-scien}

Several build-level decisions are reported here because downstream scientific claims rest on them, not because they are implementation trivia.

\textbf{Determinism.} The build is a pure function of its inputs: no randomness, stable sort orders before every deduplication, content-hashed merge identifiers, and modal-value tie-breaks resolved by deterministic ordering. Phase 2 of entity resolution is explicitly idempotent: re-running it on already-processed outputs repairs rather than duplicates canonical rows (\texttt{consolidate\_duplicate\_canonical\_orgs}). Deterministic rebuilds make the Chapter 4 oracle audit meaningful, because two independent implementations can only be compared over a reproducible substrate. They also keep the benchmark's gold programs stable across the project's lifetime.

\textbf{Validation-gated writes.} The builder runs a validation suite: primary-key uniqueness for every node table, \texttt{(ocid, award\_id)} uniqueness, referential closure for every edge endpoint, alias-target closure, value-additivity consistency, and coverage measurements. It \emph{refuses to write} the graph if any check fails, and coverage numbers are emitted as warnings into a machine-readable report (\texttt{kg/build.py}, \texttt{kg/validate.py}). The counts quoted in Section 3.2 are therefore properties the artifact is guaranteed to have, not properties it happened to have once.

\textbf{Stable composite keys.} \texttt{contract\_node\_id = contract:{ocid}:{award\_id}} names the unit of analysis, and every edge ID embeds its endpoints (\texttt{buyer:{canonical\_id}:{node\_id}}). Edges are therefore self-describing and can be deduplicated with set semantics. When one award names several raw supplier IDs that resolve to the same canonical organisation, the build collapses them into a single edge whose \texttt{raw\_id} field carries the sorted JSON list of contributing raw IDs. Multiplicity is normalised away from the graph's semantics but preserved in its provenance.

\textbf{Empty strings are a first-class hazard.} Missing values throughout the corpus arrive as empty strings, not nulls, and the build treats them accordingly. Supplier IDs that are empty or unresolvable produce no edge rather than an edge to an empty organisation. CPV edges are emitted only for non-empty codes. Derived flags (\texttt{has\_award\_signed\_date}, \texttt{has\_contract\_period}) are computed by explicit non-empty tests. This hazard escaped into computation once. In an early bridge-family oracle, an empty-string party name matched every record lacking that party, inflating every bridge sum by a constant. The dual-oracle agreement audit caught it, and Chapter 4 dissects it. It is the concrete illustration of why this chapter insists that conventions be checkable.

\section{What the knowledge graph hands to the rest of the thesis}\label{sec:what-the-knowledge-graph-hands-to-the-re}

The chapters that follow consume the substrate built here through a deliberately small contract. The \emph{grounding targets} are the typed slots and closed value spaces of Section 3.2: the planner must land on real fields, real CPV codes, and canonical organisations, and anything else is mechanically rejectable. The \emph{verifier guards} are the explicit conventions: additive-only money aggregation via \texttt{value\_is\_additive}, the flat first-party record universe, the latest-release snapshot, and the decoration-variant \texttt{IN}-filter policy. The \emph{abstention semantics} are fixed by the measured coverage bounds and by provenance: empty result sets are meaningful, expected outcomes with dedicated benchmark traps, not failure modes. None of this makes the graph "ground truth". Correctness downstream is always relative to the conventions this chapter documents, and Chapter 8 lists this relativity among the limitations. Nor does the graph claim complete coverage of UK procurement; it covers the ingested corpus only. What it does claim, supported by the artifact-verified validation gates, is that the residual ambiguity of the substrate is documented and bounded. That is the property everything downstream is built against. Chapter 4 constructs a 12,828-question benchmark by executing parameterised plans directly against this graph, and Chapter 5 builds the compile--ground--execute--verify chain over the same typed interface.
